\documentclass[a4paper, 12pt]{report}

\usepackage{amsmath}
\usepackage{amssymb}
\usepackage[margin=1in]{geometry}
\usepackage{hyperref}
\usepackage{graphicx}
\usepackage{subcaption}
\usepackage{float}
\usepackage{array}
\usepackage{booktabs}
\setcounter{secnumdepth}{3}

\renewcommand{\thesection}{\arabic{section}}
\renewcommand{\thesubsection}{\arabic{section}.\arabic{subsection}}

\setlength{\parindent}{0pt}

\begin{document}

\chapter*{Nonlinear Study}

In this chapter, the nonlinearity characteristics of Sarkic's section are studied to validate the principle of 
linear superposition and Scanlan's linearity assumption. This section is organised as follows, \\

\begin{flushleft}
\hspace{2cm}1. Motivation for Nonlinearity Study \\
\hspace{2cm}2. Simulation Parameters \\
\hspace{2cm}3. Results and Inferences \\
\end{flushleft}

\section{Motivation for Nonlinearity Study}

The base assumption of the Multi-Frequency Method is that the principle of linear superposition holds. 
This assumption is derived from Scanlan's linearity assumption, which states that for small heave and pitch amplitudes, 
the applied forced motion oscillation and the obtained aerodynamic forces are equivalent with a phase 
shift. Although the above two assumptions might appear to point to the same meaning, they can differ in 
certain cases. From Scanlan's linearity assumption, the obtained forces can be linear with respect to 
the applied motion, but the phase shift can be non-constant. 
Therefore, it becomes essential to validate the principle of linear superposition and study the 
linearity limit to ensure the Multi-Frequency method is valid and to create a framework for this method. \\

This study focuses on validating the assumption and extending the knowledge of the nonlinearity characteristics 
of Sarkic's section. Although studying the nonlinearity characteristics of a particular cross-section may make 
the developed framework specific to that cross-section, this is done so that any unexpected discrepancies 
in the results of the Multi-Frequency method can be reasonably explained. These inferences from the nonlinearity study are then generalised 
when framing the procedure for the Multi-Frequency method.\\

The study is conducted by increasing the heave and pitch motion amplitudes and then examining the obtained 
flutter derivatives, aerodynamic forces, and their phase. 
\vspace{0.5cm}

\section{Simulation Parameters}

Since the motion amplitudes are the only parameters changed, the remaining parameters derived from Chapter 2 are kept the same.
From the literature review, 5$\%$ of the section width for heave amplitude and 5 degrees for pitch are 
the linearity limit. Hence, six CFD simulations are performed by linearly increasing the heave 
amplitudes from 0.68$\%$ to 8.20$\%$ and the pitch amplitude from 1 degree to 8 degrees. 
Table~\ref{tab:nonlinear_amplitude} shows the simulation parameters for the amplitude ranges.\\

For each heave and pitch motion amplitude, aerodynamic forces are computed at 8 different reduced velocities (shown in Table~\ref{tab:Reduced velocities}). As described in 
Chapter 2, the stable time period for each of the reduced velocities is also given the table. A total of $7 \times 8 \times 2 = 112$ 
simulations are conducted. \\

\begin{table}[htbp]
    \centering
    \caption{Heave and pitch amplitude}
    \label{tab:nonlinear_amplitude}
    \begin{tabular}{|l|>{\centering\arraybackslash}p{2cm}|>{\centering\arraybackslash}p{2cm}|>{\centering\arraybackslash}p{2cm}|>{\centering\arraybackslash}p{2cm}|}
        \hline
        & \textbf{Heave Amplitude (m)} & \textbf{Heave ratio} & \textbf{Pitch (radians)} & \textbf{Pitch (Degree)} \\
        \hline
        \textbf{Amplitude 0} & 0.0025 & 0.68 & 0.017 & 0.97 \\
        \hline
        \textbf{Amplitude 1} & 0.005 & 1.36 & 0.027 & 1.55 \\
        \hline
        \textbf{Amplitude 2} & 0.01 & 2.73 & 0.054 & 3.09 \\
        \hline
        \textbf{Amplitude 3} & 0.015 & 4.09 & 0.07 & 4.01 \\
        \hline
        \textbf{Amplitude 4} & 0.0183 & 5 & 0.087 & 4.98 \\
        \hline
        \textbf{Amplitude 5} & 0.02 & 5.46 & 0.1047 & 6 \\
        \hline
        \textbf{Amplitude 6} & 0.03 & 8.20 & 0.14 & 8.02 \\
        \hline
    \end{tabular}
\end{table}

\begin{table}[htbp]
    \centering
    \caption{Reduced velocities}
    \label{tab:Reduced velocities}
    \begin{tabular}{|l|c|c|c|c|c|c|c|c|}
        \hline
        \textbf{Frequencies (Hz)} & 0.125 & 0.15 & 0.2 & 0.25 & 0.4 & 0.8 & 1.6 & 3.2 \\ \hline
        \textbf{Reduced Velocity} & 21.86 & 18.21 & 13.66 & 10.93 & 6.83 & 3.42 & 1.71 & 0.85 \\ \hline
        \textbf{Stable Time period (secs)} & 80 & 80 & 50 & 40 & 25 & 12.5 & 6.25 & 3.125 \\ \hline
    \end{tabular}
\end{table}

\section{Results and Inferences}

The following results are studied to understand the nonlinearity characteristics of the section:

\begin{flushleft}
\hspace{2cm}1. Comparison of Flutter Derivatives \\
\hspace{2cm}2. Spectral Analysis \\
\hspace{2cm}3. Aerodynamic Force and Moment Analysis \\
\hspace{2cm}4. Pressure Distribution Analysis \\
\hspace{2cm}5. Influence of Nonlinearity on Critical Velocity \\
\end{flushleft}

\subsection{Comparison of Flutter Derivatives}

Figure~\ref{fig:Combined_Flutter_Derivatives} shows the flutter derivatives obtained for 6 heave 
and pitch amplitudes. It is observed that derivatives $H_1^*$, $H_4^*$, $A_1^*$, and $A_4^*$ are equal for amplitudes up to a 
heave ratio of 5.46$\%$. When it exceeds this value, the flutter derivatives vary. This validates the thumb 
rule of a 5$\%$ linearity limit. For derivatives $H_2^*$, $H_3^*$, and $A_3^*$, the linear 
limit holds up to 3 degrees, although $H_2^*$ exhibits a slight nonlinear behaviour within that range. 
This is slightly lower than the rule of thumb of 5 degrees. \\ 
From Figure~\ref{fig:Combined_Flutter_Derivatives}, derivative $A_2^*$ significantly deviates beyond 1.55 degrees. 
It should be noted that derivative $A_2^*$ even changes its sign at Amplitude Case 3 (pitch degree = 4.01). 

\begin{figure}[htbp]  
    \centering
    \includegraphics[width=1\textwidth]{Figures/Combined_Flutter_Derivative.png}
    \caption{Combined Flutter Derivatives}
    \label{fig:Combined_Flutter_Derivatives}
\end{figure}

\subsection{Spectral Analysis}

Spectral analysis is a powerful tool to identify the presence of nonlinear terms in the obtained forces and validating scanlan's linearity assumption. 
If the spectrum of the obtained forces contains higher-order sine terms other than applied frequency, 
then nonlinear behaviour is present. In this section, the spectrum of the obtained aerodynamic lift force and moment is analysed for different amplitudes. 
Discrete Fourier Transform (DFT) is used to obtain the spectrum of the obtained forces.
In this section, plots of spectrum corresponding to the reduced velocity of 21.86 are shown, since this is the most affected by nonlinearity.
%%TODO: add complete plot of spectrum analysis. amptiude case 1,3,5 are missing

From Figure~\ref{fig:Frequency_Domain_Analysis_Heave} for all heave ampltiude range, the spectrum of the obtained both aerodynamic forces and moments are clean, 
indicating that the obtained forces are linear with respect to the applied motion.

From Figure~\ref{fig:Frequency_Domain_Analysis_Pitch}, for pitch amplitude up to 5 degrees for pitch, the spectrum of the obtained both aerodynamic forces and moments are clean, 
indicating that the obtained forces are linear with respect to the applied motion. Above these limits, significant noise is present, indicating the 
presence of higher-order sine terms and, therefore nonlinearity starts to appear. \\

On comparing Figure~\ref{fig:Combined_Flutter_Derivatives} and Figure~\ref{fig:Frequency_Domain_Analysis_Pitch}, 
it can be observed that the nonlinearity in flutter derivatives $A_2^*$ and $H_2^*$ starts to deviate at a lower pitch degree than the nonlinearity in the obtained spectrum analysis.
Hence this deviation in flutter derivatives is not due to nonlinearity in the obtained forces. This validdates the rule of thumb of 5 degrees for pitch motion, 
but it also raises the question of why the flutter derivatives show nonlinearity at a lower pitch degree than the obtained forces.

\begin{figure}[htbp]  
    \centering
    \begin{subfigure}[b]{0.45\textwidth}
        \centering
        \includegraphics[width=\textwidth]{Figures/Frequency_Domain_Analysis/Heave/DFT_LiftForce_Amplitude_Case_0.png}
        \caption{}
        \label{fig:Heave_DFT_Fy_amplitude_case_0}
    \end{subfigure}
    \hfill
    \begin{subfigure}[b]{0.45\textwidth}
        \centering
        \includegraphics[width=\textwidth]{Figures/Frequency_Domain_Analysis/Heave/DFT_Moment_Amplitude_Case_0.png}
        \caption{}
        \label{fig:Heave_DFT_moment_amplitude_case_0}
    \end{subfigure}
    
    \begin{subfigure}[b]{0.45\textwidth}
        \centering
        \includegraphics[width=\textwidth]{Figures/Frequency_Domain_Analysis/Heave/DFT_LiftForce_Amplitude_Case_1.png}
        \caption{}
        \label{fig:Heave_DFT_Fy_amplitude_case_2}
    \end{subfigure}
    \hfill
    \begin{subfigure}[b]{0.45\textwidth}
        \centering
        \includegraphics[width=\textwidth]{Figures/Frequency_Domain_Analysis/Heave/DFT_Moment_Amplitude_Case_1.png}
        \caption{}
        \label{fig:Heave_DFT_moment_amplitude_case_2}
    \end{subfigure}
    
    \begin{subfigure}[b]{0.45\textwidth}
        \centering
        \includegraphics[width=\textwidth]{Figures/Frequency_Domain_Analysis/Heave/DFT_LiftForce_Amplitude_Case_2.png}
        \caption{}
        \label{fig:Heave_DFT_Fy_amplitude_case_4}
    \end{subfigure}
    \hfill
    \begin{subfigure}[b]{0.45\textwidth}
        \centering
        \includegraphics[width=\textwidth]{Figures/Frequency_Domain_Analysis/Heave/DFT_Moment_Amplitude_Case_2.png}
        \caption{}
        \label{fig:Heave_DFT_moment_amplitude_case_4}
    \end{subfigure}
    
    \begin{subfigure}[b]{0.45\textwidth}
        \centering
        \includegraphics[width=\textwidth]{Figures/Frequency_Domain_Analysis/Heave/DFT_LiftForce_Amplitude_Case_3.png}
        \caption{}
        \label{fig:Heave_DFT_Fy_amplitude_case_6}
    \end{subfigure}
    \hfill
    \begin{subfigure}[b]{0.45\textwidth}
        \centering
        \includegraphics[width=\textwidth]{Figures/Frequency_Domain_Analysis/Heave/DFT_Moment_Amplitude_Case_3.png}
        \caption{}
        \label{fig:Heave_DFT_moment_amplitude_case_6}
    \end{subfigure}

    \begin{subfigure}[b]{0.45\textwidth}
        \centering
        \includegraphics[width=\textwidth]{Figures/Frequency_Domain_Analysis/Heave/DFT_LiftForce_Amplitude_Case_4.png}
        \caption{}
        \label{fig:Heave_DFT_Fy_amplitude_case_6}
    \end{subfigure}
    \hfill
    \begin{subfigure}[b]{0.45\textwidth}
        \centering
        \includegraphics[width=\textwidth]{Figures/Frequency_Domain_Analysis/Heave/DFT_Moment_Amplitude_Case_4.png}
        \caption{}
        \label{fig:Heave_DFT_moment_amplitude_case_6}
    \end{subfigure}

    \begin{subfigure}[b]{0.45\textwidth}
        \centering
        \includegraphics[width=\textwidth]{Figures/Frequency_Domain_Analysis/Heave/DFT_LiftForce_Amplitude_Case_5.png}
        \caption{}
        \label{fig:Heave_DFT_Fy_amplitude_case_6}
    \end{subfigure}
    \hfill
    \begin{subfigure}[b]{0.45\textwidth}
        \centering
        \includegraphics[width=\textwidth]{Figures/Frequency_Domain_Analysis/Heave/DFT_Moment_Amplitude_Case_5.png}
        \caption{}
        \label{fig:Heave_DFT_moment_amplitude_case_6}
    \end{subfigure}    

    \begin{subfigure}[b]{0.45\textwidth}
        \centering
        \includegraphics[width=\textwidth]{Figures/Frequency_Domain_Analysis/Heave/DFT_LiftForce_Amplitude_Case_6.png}
        \caption{}
        \label{fig:Heave_DFT_Fy_amplitude_case_6}
    \end{subfigure}
    \hfill
    \begin{subfigure}[b]{0.45\textwidth}
        \centering
        \includegraphics[width=\textwidth]{Figures/Frequency_Domain_Analysis/Heave/DFT_Moment_Amplitude_Case_6.png}
        \caption{}
        \label{fig:Heave_DFT_moment_amplitude_case_6}
    \end{subfigure}

    \caption{Frequency Domain Analysis Heave}
    \label{fig:Frequency_Domain_Analysis_Heave}
\end{figure}

\begin{figure}[htbp]  
    \centering
    \begin{subfigure}[b]{0.45\textwidth}
        \centering
        \includegraphics[width=\textwidth]{Figures/Frequency_Domain_Analysis/Pitch/DFT_LiftForce_Amplitude_Case_0.png}
        \caption{}
        \label{fig:Pitch_DFT_Fy_amplitude_case_0}
    \end{subfigure}
    \hfill
    \begin{subfigure}[b]{0.45\textwidth}
        \centering
        \includegraphics[width=\textwidth]{Figures/Frequency_Domain_Analysis/Pitch/DFT_Moment_Amplitude_Case_0.png}
        \caption{}
        \label{fig:Pitch_DFT_moment_amplitude_case_0}
    \end{subfigure}
    
    \begin{subfigure}[b]{0.45\textwidth}
        \centering
        \includegraphics[width=\textwidth]{Figures/Frequency_Domain_Analysis/Pitch/DFT_LiftForce_Amplitude_Case_1.png}
        \caption{}
        \label{fig:Pitch_DFT_Fy_amplitude_case_2}
    \end{subfigure}
    \hfill
    \begin{subfigure}[b]{0.45\textwidth}
        \centering
        \includegraphics[width=\textwidth]{Figures/Frequency_Domain_Analysis/Pitch/DFT_Moment_Amplitude_Case_1.png}
        \caption{}
        \label{fig:Pitch_DFT_moment_amplitude_case_2}
    \end{subfigure}
    
    \begin{subfigure}[b]{0.45\textwidth}
        \centering
        \includegraphics[width=\textwidth]{Figures/Frequency_Domain_Analysis/Pitch/DFT_LiftForce_Amplitude_Case_2.png}
        \caption{}
        \label{fig:Pitch_DFT_Fy_amplitude_case_4}
    \end{subfigure}
    \hfill
    \begin{subfigure}[b]{0.45\textwidth}
        \centering
        \includegraphics[width=\textwidth]{Figures/Frequency_Domain_Analysis/Pitch/DFT_Moment_Amplitude_Case_2.png}
        \caption{}
        \label{fig:Pitch_DFT_moment_amplitude_case_4}
    \end{subfigure}
    
    \begin{subfigure}[b]{0.45\textwidth}
        \centering
        \includegraphics[width=\textwidth]{Figures/Frequency_Domain_Analysis/Pitch/DFT_LiftForce_Amplitude_Case_3.png}
        \caption{}
        \label{fig:Pitch_DFT_Fy_amplitude_case_6}
    \end{subfigure}
    \hfill
    \begin{subfigure}[b]{0.45\textwidth}
        \centering
        \includegraphics[width=\textwidth]{Figures/Frequency_Domain_Analysis/Pitch/DFT_Moment_Amplitude_Case_3.png}
        \caption{}
        \label{fig:Pitch_DFT_moment_amplitude_case_6}
    \end{subfigure}

    \begin{subfigure}[b]{0.45\textwidth}
        \centering
        \includegraphics[width=\textwidth]{Figures/Frequency_Domain_Analysis/Pitch/DFT_LiftForce_Amplitude_Case_4.png}
        \caption{}
        \label{fig:Pitch_DFT_Fy_amplitude_case_6}
    \end{subfigure}
    \hfill
    \begin{subfigure}[b]{0.45\textwidth}
        \centering
        \includegraphics[width=\textwidth]{Figures/Frequency_Domain_Analysis/Pitch/DFT_Moment_Amplitude_Case_4.png}
        \caption{}
        \label{fig:Pitch_DFT_moment_amplitude_case_6}
    \end{subfigure}

    \begin{subfigure}[b]{0.45\textwidth}
        \centering
        \includegraphics[width=\textwidth]{Figures/Frequency_Domain_Analysis/Pitch/DFT_LiftForce_Amplitude_Case_5.png}
        \caption{}
        \label{fig:Pitch_DFT_Fy_amplitude_case_6}
    \end{subfigure}
    \hfill
    \begin{subfigure}[b]{0.45\textwidth}
        \centering
        \includegraphics[width=\textwidth]{Figures/Frequency_Domain_Analysis/Pitch/DFT_Moment_Amplitude_Case_5.png}
        \caption{}
        \label{fig:Pitch_DFT_moment_amplitude_case_6}
    \end{subfigure}    

        \begin{subfigure}[b]{0.45\textwidth}
        \centering
        \includegraphics[width=\textwidth]{Figures/Frequency_Domain_Analysis/Pitch/DFT_LiftForce_Amplitude_Case_6.png}
        \caption{}
        \label{fig:Pitch_DFT_Fy_amplitude_case_6}
    \end{subfigure}
    \hfill
    \begin{subfigure}[b]{0.45\textwidth}
        \centering
        \includegraphics[width=\textwidth]{Figures/Frequency_Domain_Analysis/Pitch/DFT_Moment_Amplitude_Case_6.png}
        \caption{}
        \label{fig:Pitch_DFT_moment_amplitude_case_6}
    \end{subfigure}

    \caption{Frequency Domain Analysis Pitch}
    \label{fig:Frequency_Domain_Analysis_Pitch}
\end{figure}

\subsection{Aerodynamic Force and Moment Analysis}

In this section, the obtained aerodynamic forces and moments are compared against the applied motion 
amplitude to understand the nonlinearity characteristics of forces and moments. This is also 
conducted for phase of the aerodynamic forces and moments.

Figure~\ref{fig:Displacement_Vs_Force} shows plots of applied amplitude vs aerodynamic force and moment, and 
applied amplitude vs phase of aerodynamic force and moment, for both heave and pitch motion. These are plotted for 
a reduced velocity of 21.86. It can be observed that for heave motion, the lift force and moment amplitude 
exhibit linear behaviour with respect to the applied displacement. Furthermore, the phase of the force and 
moment is constant. For pitch motion, the lift force and moment shows linear behaviour upto to the pitch degree of 5
(the moment exhibits very little nonlinearity in this region). This is again consistent with the rule of thumb of 5 degrees and 
spectral analysis and therefore validating the scanlan's linearity assumption. \\

However, the phase is not constant for either the lift force or the moment after 3 degrees. To have a 
comparative study of phase with flutter derivatives, first their relaton is once again stated, then comparitive study is shown. 
Derivatives $H_1^*$, $A_1^*$, $H_4^*$, and $A_4^*$ are associated with the heave motion aerodynamic 
lift, moment, and their phase. Derivatives $H_2^*$, $A_2^*$, $H_3^*$, and $A_3^*$ are associated with the pitch motion aerodynamic 
lift, moment, and their phase. $H_2^*$ and $A_2^*$ are associated with the phase of the pitch motion aerodynamic 
lift and moment. \\

On comparing Figure~\ref{fig:Displacement_Vs_Force} and Figure~\ref{fig:Combined_Flutter_Derivatives}, 
it is clear that the nonlinearity in flutter derivatives $A_2^*$ and $H_2^*$ is due to the phase change.
Further, on examining~\ref{fig:Displacement_Vs_Force}{h} the phase of the moment (related to derivative $A_2^*$), it changes 
from positive to negative at 3 degrees. And it can be observed in Figure~\ref{fig:Combined_Flutter_Derivatives} 
at 3 degrees derivative $A_2^*$ changes its sign. This confirms the sign change in $A_2^*$ is caused by the phase change. \\

It is therefore concluded that the nonlinearity in derivatives $A_2^*$ and $H_2^*$ is not due to nonlinearity in the obtained forces, 
but due to the phase change.
This is reason the spectrum of the obtained forces is clean: the amplitude of the force is linear, but the phase is not.
A similar behaviour was observed by --- and ---. \\

This observation points to the earlier discussion that Scanlan's linearity assumption is not equivalent to the principle of linear superposition, 
since the obtained forces can be linear with respect to the applied motion, but the phase shift can be non-constant. \\

Since the plot shown before was only for the redcued velocity of 21.86, it is necessary to analyse the results for other reduced 
velocities before deducing the above conclusion. The results for all reduced velocities are compiled in Figure~\ref{fig:Compiled_Displacement_Vs_Force}. 
The x-axis is same as before, the ampltiude case, but the y-axis is modified.
This y-axis represents aerodynamic force divided by the applied motion amplitude and normalised with respect to the first amplitude case. 
The aerodynamic force is normalised with respect to the first amplitude case because the different reduced velocities will yield different constants 
and normalisation allows for a better comparison of the nonlinearity characteristics across different reduced velocities.
In simple terms, if the aerodynamic forces are linear, the plot should be a constant horizontal line. If the phase is constant, 
the plot should similarly be a horizontal line.\\
    
It can be observed that for heave motion, lift and moment, and their respective phases remain constant up to heave ratio of 5.46 
for reduced velocities above 1.71. 
The constancy of the remaining cases replicates the behaviour observed earlier for the reduced velocity of 21.86
(Figure~\ref{fig:Displacement_Vs_Force}). For pitch motion, the lift force and moment exhibit slight nonlinearity, 
as seen for the reduced velocity of 21.86 (Figure~\ref{fig:Displacement_Vs_Force}).  
A similar behaviour is observed for the phase, which is also not constant. 

The reduced velocities of 1.71 (frequency = 1.6 Hz) and 0.85 (frequency = 3.2 Hz) alone shows significant deviation 
in both heave and pitch motions.  
As they are affected by the vortex shedding effect, capturing these derivatives is difficult, and therefore they 
are more affected by noise.
Another important observation is that the nonlinearity in phase increases with increasing reduced velocity (decreasing frequency).
This further validates the behaviour observed in Figure~\ref{fig:Combined_Flutter_Derivatives}, where the nonlinearity in $A_2^*$ and $H_2^*$ is lower at lower reduced velocities. \\

To conclude, although Scanlan's assumption of linearity holds for the 
obtained forces, aerodynamic derivatives may still exhibit nonlinearity. And the essential note is that this 
nonlinear behaviour is not detectable with spectrum analysis. \\

% Nonlinearity Study
\begin{figure}[htbp]  
    \centering
    \begin{subfigure}[b]{0.45\textwidth}
        \centering
        \includegraphics[width=\textwidth]{Figures/Displacement_Vs_Force/Heave/Fy_vs_Amplitude.png}
        \caption{}
        \label{fig:Heave_Fy_vs_Amplitude}
    \end{subfigure}
    \hfill
    \begin{subfigure}[b]{0.45\textwidth}
        \centering
        \includegraphics[width=\textwidth]{Figures/Displacement_Vs_Force/Pitch/Fy_vs_Amplitude.png}
        \caption{}
        \label{fig:Pitch_Fy_vs_Amplitude}
    \end{subfigure}
    
    \begin{subfigure}[b]{0.45\textwidth}
        \centering
        \includegraphics[width=\textwidth]{Figures/Displacement_Vs_Force/Heave/PhaseFy_vs_Amplitude.png}
        \caption{}
        \label{fig:Heave_PhaseFy_vs_Amplitude}
    \end{subfigure}
    \hfill
    \begin{subfigure}[b]{0.45\textwidth}
        \centering
        \includegraphics[width=\textwidth]{Figures/Displacement_Vs_Force/Pitch/PhaseFy_vs_Amplitude.png}
        \caption{}
        \label{fig:Pitch_PhaseFy_vs_Amplitude}
    \end{subfigure}
    
    \begin{subfigure}[b]{0.45\textwidth}
        \centering
        \includegraphics[width=\textwidth]{Figures/Displacement_Vs_Force/Heave/Moment_vs_Amplitude.png}
        \caption{}
        \label{fig:Heave_Moment_vs_Amplitude}
    \end{subfigure}
    \hfill
    \begin{subfigure}[b]{0.45\textwidth}
        \centering
        \includegraphics[width=\textwidth]{Figures/Displacement_Vs_Force/Pitch/Moment_vs_Amplitude.png}
        \caption{}
        \label{fig:Pitch_Moment_vs_Amplitude}
    \end{subfigure}
    
    \begin{subfigure}[b]{0.45\textwidth}
        \centering
        \includegraphics[width=\textwidth]{Figures/Displacement_Vs_Force/Heave/PhaseMx_vs_Amplitude.png}
        \caption{}
        \label{fig:Heave_PhaseMx_vs_Amplitude}
    \end{subfigure}
    \hfill
    \begin{subfigure}[b]{0.45\textwidth}
        \centering
        \includegraphics[width=\textwidth]{Figures/Displacement_Vs_Force/Pitch/PhaseMx_vs_Amplitude.png}
        \caption{}
        \label{fig:Pitch_PhaseMx_vs_Amplitude}
    \end{subfigure}
    
    \caption{Displacement Vs Force}
    \label{fig:Displacement_Vs_Force}
\end{figure}

% Compressed Nonlinearity Study
\begin{figure}[htbp]  
    \centering
    \begin{subfigure}[b]{0.45\textwidth}
        \centering
        \includegraphics[width=\textwidth]{Figures/Compressed_Displacement_Vs_Force/Heave/NormalizedFy_vs_Amplitude_Case_RelativeToFirst.png}
        \caption{}
        \label{fig:Compiled_Heave_Fy_vs_Amplitude}
    \end{subfigure}
    \hfill
    \begin{subfigure}[b]{0.45\textwidth}
        \centering
        \includegraphics[width=\textwidth]{Figures/Compressed_Displacement_Vs_Force/Pitch/NormalizedFy_vs_Amplitude_Case_RelativeToFirst.png}
        \caption{}
        \label{fig:Compiled_Pitch_Fy_vs_Amplitude}
    \end{subfigure}
    
    \begin{subfigure}[b]{0.45\textwidth}
        \centering
        \includegraphics[width=\textwidth]{Figures/Compressed_Displacement_Vs_Force/Heave/PhaseFy_vs_Amplitude_Case_RelativeToFirst.png}
        \caption{}
        \label{fig:Compiled_Heave_PhaseFy_vs_Amplitude}
    \end{subfigure}
    \hfill
    \begin{subfigure}[b]{0.45\textwidth}
        \centering
        \includegraphics[width=\textwidth]{Figures/Compressed_Displacement_Vs_Force/Pitch/PhaseFy_vs_Amplitude_Case_RelativeToFirst.png}
        \caption{}
        \label{fig:Compiled_Pitch_PhaseFy_vs_Amplitude}
    \end{subfigure}
    
    \begin{subfigure}[b]{0.45\textwidth}
        \centering
        \includegraphics[width=\textwidth]{Figures/Compressed_Displacement_Vs_Force/Heave/NormalizedMoment_vs_Amplitude_Case_RelativeToFirst.png}
        \caption{}
        \label{fig:Compiled_Heave_Moment_vs_Amplitude}
    \end{subfigure}
    \hfill
    \begin{subfigure}[b]{0.45\textwidth}
        \centering
        \includegraphics[width=\textwidth]{Figures/Compressed_Displacement_Vs_Force/Pitch/NormalizedMoment_vs_Amplitude_Case_RelativeToFirst.png}
        \caption{}
        \label{fig:Compiled_Pitch_Moment_vs_Amplitude}
    \end{subfigure}
    
    \begin{subfigure}[b]{0.45\textwidth}
        \centering
        \includegraphics[width=\textwidth]{Figures/Compressed_Displacement_Vs_Force/Heave/PhaseMx_vs_Amplitude_Case_RelativeToFirst.png}
        \caption{}
        \label{fig:Compiled_Heave_PhaseMx_vs_Amplitude}
    \end{subfigure}
    \hfill
    \begin{subfigure}[b]{0.45\textwidth}
        \centering
        \includegraphics[width=\textwidth]{Figures/Compressed_Displacement_Vs_Force/Pitch/PhaseMx_vs_Amplitude_Case_RelativeToFirst.png}
        \caption{}
        \label{fig:Compiled_Pitch_PhaseMx_vs_Amplitude}
    \end{subfigure}
    
    \caption{Compiled Displacement Vs Force}
    \label{fig:Compiled_Displacement_Vs_Force}
\end{figure}


\subsection{Pressure Distribution Analysis}
% TODO: add pressure distribution analysis - pictures of pressure distribution on cross section over different time steps 
As mentioned, a similar behaviour was observed by --- and ---. However, the observed phase deviation in those studies occurred at a higher pitch degree. 
This raises the question of whether these results stem from the same behaviour or from different underlying reasons. 
The reason for the observed behaviour in --- is explained in the paper itself.
It is attributed to the reattachment point of the flow. The flow separates at the leading edge and reattaches 
at the trailing edge for lower pitch degrees, but as the pitch degree increases, the reattachment point moves towards the leading edge. 
This shift in the reattachment point causes the phase change, which is the reason for the observed nonlinearity. \\

To confirm whether this is also the reason in the present case, the pressure distribution is plotted for different 
pitch degrees. It is observed that the reattachment point moves towards the leading edge as the pitch 
degree increases, which confirms that the nonlinearity in this case is also attributable to the 
change in reattachment point. \\

Finally, it should be noted that the linearity limit for pitch motion is considerably lower than the rule of thumb of 5 degrees.
A preliminary comparative study of three sections suggests that the breadth-to-depth ratio may play an important 
role in this behaviour. CFD modelling and turbulence modelling are unlikely to be the cause, since both referenced papers are based on wind tunnel experiments. \\



\subsection{Influence of Nonlinearity on Critical Velocity}

Lastly, the influence of this nonlinearity on the critical velocity is studied. 
For the calculation of critical velocity, structural parameters are required. 
Five bridge samples are considered, and their structural properties are summarised in Table~\ref{tab:bridge_configs}.

\begin{table}[h!]
\centering
\caption{Structural parameters of bridge samples used for critical velocity analysis}
\label{tab:bridge_configs}
\begin{tabular}{lrrrrrrr}
\hline
\textbf{Bridge} & \textbf{$B$ (m)} & \textbf{$m$ (kg/m)} & \textbf{$\zeta_h$} & \textbf{$f_h$ (Hz)} & \textbf{$I$ (kg$\cdot$m)} & \textbf{$\zeta_\alpha$} & \textbf{$f_\alpha$ (Hz)} \\
\hline
Diana (2020)       & 31.0  & 22740  & 0.003 & 0.100  & $2.47 \times 10^6$  & 0.003 & 0.278 \\
Great Belt         & 18.3  & 12800  & 0.003 & 0.1408 & $4.28 \times 10^5$  & 0.003 & 0.360 \\
Great Belt (v2)    & 18.3  & 23687  & 0.003 & 0.097  & $2.50 \times 10^6$  & 0.003 & 0.270 \\
Humen              & 35.6  & 18300  & 0.005 & 0.1117 & $2.09 \times 10^6$  & 0.005 & 0.361 \\
Tacoma Narrows     & 11.9  & 4250   & 0.005 & 0.130  & $1.78 \times 10^5$  & 0.005 & 0.200 \\
\hline
\end{tabular}
\end{table}

\begin{figure}[htbp]  
    \centering
    \includegraphics[width=1\textwidth]{Figures/CriticalVelocity_vs_Number_of_Combinations.png}
    \caption{Critical Velocity vs Number of Combinations}
    \label{fig:CriticalVelocity_vs_Number_of_Combinations}
\end{figure}

\section{Observations and Discussion}

\begin{enumerate}
    \item The linearity limit for heave motion agrees with the rule of thumb; therefore, there should not be any issue with the Multi-Frequency method as long as the amplitude remains within this limit.
    
    \item Scanlan's assumptions and the principle of superposition are not equivalent and may differ; this is most clearly seen with pitch motion.
    
    \item Therefore, when validating the Multi-Frequency method, it is not necessary to obtain accurate derivatives for $A_2^*$ and $H_2^*$.
    
    \item Although the nonlinearity in $A_2^*$ and $H_2^*$ has been established in multiple research papers, the important finding here is that the present model exhibits nonlinearity at a much lower pitch degree. While the reason for this behaviour is linked to the reattachment point, an important question remains: what governs the onset of this behaviour, and why do different studies report different nonlinearity limits? This may be related to sectional properties rather than CFD modelling choices.
\end{enumerate}

\end{document}